\section{Implications and Conclusion}\label{sec:conclusion}

\subsection{Scope}\label{sec:scope}

Four conditions define the class of cases: (i) \emph{security dependence}, protection the holder cannot replace on the relevant horizon ($S_k>0$); (ii) \emph{a partly conditional guarantee}, some of which the patron can stake to withhold after exit ($\kappa S_k+d_p>0$; \cref{prop:conditionality}); (iii) \emph{a rival that bids} ($\omega>0$); and (iv) \emph{a nested chokepoint}, partly inside the patron's jurisdiction while the holder's licence still decides exit ($0<\theta<1$). The Netherlands, Japan and South Korea meet the first and fourth. On the second, the May 2026 posture order against Germany (\cref{sec:background}) suggests that $d_p$ is positive; the conditionality of the core guarantee is largely untested (\cref{sec:record}). The third holds only weakly: we found no Chinese offer for any holder's exit, one reason we think $\omega$ is small. Taiwan (threatened by China, not courted) and the Gulf partners (no treaty guarantee) fall outside.

\subsection{Implications for Europe}\label{sec:implications}

We place Europe at a nonpositive premium with little rupture risk: The Hague has received revocable statuses, which we do not code as transfers (\cref{sec:coding}), and has made no demand whose rejection risked realignment (\cref{sec:record}). Whatever settles Europe's price, its surplus over exit bounds what Washington's conditional threats can extract (\cref{prop:slack}).

\paragraph{Reduce the conditional part of dependence.} Against the NATO target of 5 percent of GDP by 2035 \citep{nato2025hague}, the model counts a euro only insofar as it replaces something Washington could withhold after exit: substitutes for posture and enablers more than duplicates of the statutory core (\cref{sec:security}), and AI capacity Washington cannot withhold, which lowers $d_{ai}$. While $Q<0$ a gain in premium removes slack one for one and brings no transfer (\cref{cor:rationing}). Where assurance is cheap at the margin, Washington answers with a more conditional guarantee and pays nothing (\cref{prop:conditionality}(c)(i)).

\paragraph{Keep concessions reversible together.} Because exit would undo Europe's reversible concessions at once, they share one budget, which protects Europe against an accumulation of demands. A concession that Washington can enforce with a penalty of its own, such as a snapback tariff, escapes that budget and costs Europe more (\cref{prop:linkage}). Europe gains from resisting issue-specific enforcement clauses even where it accepts the concession itself.

\paragraph{Shrink the jurisdictional source of slack.} Washington reaches the firm through US-origin inputs and US persons (\cref{sec:background}); on our reading that shared jurisdiction can keep the premium negative, and so create slack, without any security term. Moving that content and servicing out of US jurisdiction lowers $\theta$, which raises $\omega$ only while China's cap does not bind (\cref{eq:omega}); the surer gain is a lower share $h$ of servicing that Washington can compel (\cref{sec:holdup}). In May 2026 The Hague opposed the MATCH Act's extraterritorial reach, arguing that ``every country is responsible for its own export control legislation'' \citep[our translation]{tk2026match}, which we read as a defence of who holds the key.

\paragraph{Interrupt servicing only under a rule that cannot be suspended.} Under the Union's timing, interruption helps only if both sides expect a rejected demand to be met by one (\cref{prop:range}), and each countermeasure the Union armed against Washington was suspended once Washington moved (\cref{sec:council,sec:holdup}). A credible option would need a legal basis, adopted before any dispute, that no EU body could suspend on a US offer, and Germany, whose Zeiss optics and Trumpf lasers cannot be commandeered abroad (\cref{sec:holdup}).

\paragraph{Expect a preventive demand; bind with care.} Under rupture risk the motive to demand control needs no positive premium, only a surplus that a partial transfer would remove (\cref{rem:controlsettle}), and autonomy investments Europe cannot commit to forgo strengthen it (\cref{cor:control}(b)). Keeping the key in Europe, or binding as Delegated Regulation 2025/2003 \citep{eu2025dualuse} did partway (\cref{sec:bind}), deters a demand for control only if Washington values control little for its own sake (\cref{rem:bind}). US demands shifting from alignment toward ownership or a licensing veto would indicate a high $v_c$, and binding would then have raised the bite and pre-empted nothing.

\paragraph{Chinese counter-threats.} If Washington has a divisible sunk demand China cannot reach, as we read the spending target and purchases of US equipment (\cref{sec:slackprinciple}), de-risking, the aim of the Critical Raw Materials Act \citep{eu2024crma}, leaves Europe's loss at $B(J^c)$ (\cref{prop:counter}(a)). It pays mainly where China retaliates against Europe's own acts, as in Nexperia (\cref{sec:record}), and it weakens China's protection of European ownership (\cref{cor:control}(d)).

\subsection{Limitations}\label{sec:limits}

Complete information about conditionality is the strongest remaining assumption. \Cref{prop:refusal} lets Washington be unsure what a demand costs Europe, which produces refusals and executed threats. If $\kappa$ were private, Washington would gain from overstating it, European beliefs would set $Q$, and the 2025 National Security Strategy would be partly cheap talk. Refusals and executed counter-threats could then come from misjudgment, the reading of Nexperia outside the model (\cref{sec:record}), and conditionality could exist only in the holder's beliefs, as Bonn's fears in 1975--77 suggest (\cref{sec:cases}). The great powers are unitary, although Congress can set stakes the executive would not, as the MATCH Act's 150-day deadline would (\cref{sec:background}). With one rival and holders bargaining alone, the model is silent on whether a coalition of the Netherlands, Japan and South Korea would help.

Within these limits, the result adds a condition to the theory of weaponized interdependence. \citet{farrell2019weaponized} require jurisdiction over a hub and the domestic institutions to use it. When a hub's holder bargains with a state that depends on it, a third condition applies: the holder's exit premium must be positive, unless it has agenda power under rupture risk (\cref{cor:master}(c)) or can impose hold-up (\cref{prop:pain}). The Netherlands licenses ASML's exports to China, so it has the institutions in that direction; turning the tools against Washington takes a Union decision (\cref{sec:council}). We think both substitutes for a positive premium are weak in its case. Whether it meets the third condition then turns on China's credible promises and, if Washington's threat to withhold frontier AI is credible, on the AI instrument, which on our orders of magnitude weighs more than the conditional part of the guarantee (\cref{sec:record}).
